\subsection{从位移矩阵元到条件光子数}
前面的 Laguerre 闭式给出全阶单模散射。为把每个阶乘和相位直接对应到跃迁路径，
令 $\zeta=-g_{\mathrm Q}^*$、$\mu=|\zeta|^2$，在固定总交换数的链上使用普通振子位移 $D(\zeta)$。
\begin{proposition}[有限和形式的位移振幅]
对 $m,n\ge0$，
\begin{equation}
 D_{mn}(\zeta)=\ee^{-\mu/2}\sqrt{m!n!}
 \sum_{k=0}^{\min(m,n)}
 \frac{\zeta^{m-k}(-\zeta^*)^{n-k}}{k!(m-k)!(n-k)!}.
 \label{eq:feqo-displacement-matrix}
\end{equation}
入射 $|0,n\rangle$ 的输出为 $\sum_mD_{mn}|n-m,m\rangle$。
\end{proposition}
\begin{proof}
因 $[\zeta a^\dagger,-\zeta^*a]=\mu I$，正规排序给
$D(\zeta)=\ee^{-\mu/2}\ee^{\zeta a^\dagger}\ee^{-\zeta^*a}$。
右侧指数从 $|n\rangle$ 降到 $|k\rangle$，系数为
$(-\zeta^*)^{n-k}\sqrt{n!/k!}/(n-k)!$；左侧再升到 $|m\rangle$，
系数为 $\zeta^{m-k}\sqrt{m!/k!}/(m-k)!$。
对允许的 $0\le k\le\min(m,n)$ 求和即得。
位移算符酉性还给 $\sum_m|D_{mn}|^2=1$，故概率由全部末态承接，不能任意归一化有限的三个边带。
\end{proof}

\begin{exbox}[单光子输入的中心峰相消]
取电子单格与光子数态 $n=1$ 入射，求中心峰和相邻两条边带的全阶振幅，判断中心峰是否总保留非零权重。

\textbf{解。} 在\eqnrefs{eq:feqo-displacement-matrix}中代入 $n=1$，有限和至多含两条路径：
\begin{equation}
 D_{0,1}=-\zeta^*\ee^{-\mu/2},\qquad
 D_{1,1}=(1-\mu)\ee^{-\mu/2},\qquad
 D_{2,1}=\frac{\zeta(2-\mu)}{\sqrt2}\ee^{-\mu/2}.
 \label{eq:feqo-single-photon-exact-amplitudes}
\end{equation}
所以 $P_{+1}=\mu\ee^{-\mu}$、$P_0=(1-\mu)^2\ee^{-\mu}$、
$P_{-1}=\mu(2-\mu)^2\ee^{-\mu}/2$。
在 $\mu=1$ 时无交换峰为零：这是振幅中不交换路径与返回路径相消，
不是光子不存在；概率仍由其他所有边带承担。
\end{exbox}

\begin{definition}[电子结果对应的条件光态]
若 $|\Psi\rangle=\sum_{\ell,m}C_{\ell m}|\ell,m\rangle$，理想电子结果 $\ell$ 的概率与条件态为
\begin{equation}
 p_\ell=\sum_m|C_{\ell m}|^2,\qquad
 |\gamma_\ell\rangle=p_\ell^{-1/2}\sum_mC_{\ell m}|m\rangle,\qquad p_\ell>0.
 \label{eq:feqo-conditional-photon-state}
\end{equation}
条件光子数分布为
\begin{equation}
 P(m|\ell)=|C_{\ell m}|^2/p_\ell.
 \label{eq:feqo-conditional-photon-number}
\end{equation}
零概率结果不定义归一化条件态；有限能窗和混合输入应使用测量章的条件量子仪器。
\end{definition}

\begin{exbox}[相干输入的精确条件分布及经典极限的限制]
取电子单格入射，光场为相干态 $|\alpha\rangle$。给定电子末阶数 $\ell$，求条件光子数分布，并判断它何时接近输入的 Poisson 分布。

\textbf{解。} 输入相干态的数态系数为 $c_n=\ee^{-|\alpha|^2/2}\alpha^n/\sqrt{n!}$。
守恒量要求同一末态 $(\ell,m)$ 只能由初态光子数 $m+\ell$ 到达，故
\begin{equation}
 C_{\ell m}=c_{m+\ell}D_{m,m+\ell}(\zeta),\qquad
 (\rho_e)_{\ell\ell'}=\sum_mC_{\ell m}C_{\ell'm}^*.
 \label{eq:feqo-coherent-reduced-exact}
\end{equation}
负指标 $c_n$ 取零。于是
\begin{equation}
 P(m|\ell)=\frac{|c_{m+\ell}D_{m,m+\ell}(\zeta)|^2}
 {\sum_{m'}|c_{m'+\ell}D_{m',m'+\ell}(\zeta)|^2}.
 \label{eq:feqo-pinem-conditional}
\end{equation}
一般有限耦合下，此分布没有理由是均值 $|\alpha|^2-\ell$ 的 Poisson 分布：
电子后选择还乘上依赖 $m$ 的位移矩阵元。
经典缩放 $g_{\mathrm Q}\to0$、$|\alpha|\to\infty$ 且 $g_{\mathrm Q}\alpha=\beta$ 固定时，
对具有非零极限概率的边带，条件光态趋近输入的宏观相干场，条件数分布与输入 Poisson 数分布趋近。
但若某条边带接近零概率，相同的无条件误差不能直接控制其归一化条件态，
误差会被 $1/p_\ell$ 放大。真空输入则完全不同：结果 $\ell=-m$ 精确预告光子数态 $|m\rangle$，
其 $g^{(2)}(0)=1-1/m$ 仅在 $m\ge1$ 时有定义；$m=0$ 的分母为零，不能称其 $g^{(2)}=0$。
\end{exbox}

\begin{proposition}[弱量子交换读取的相干与数统计]
对初始电子 $|0\rangle$，展开 $S=I+g_{\mathrm Q}b^\dagger a-g_{\mathrm Q}^*ba^\dagger+O(|g_{\mathrm Q}|^2)$。
其中电子增益振幅为 $g_{\mathrm Q}a$，电子损失振幅为 $-g_{\mathrm Q}^*a^\dagger$，因此
\begin{equation}
 P_{+1}=|g_{\mathrm Q}|^2\bar n+O(|g_{\mathrm Q}|^4),\quad
 P_{-1}=|g_{\mathrm Q}|^2(\bar n+1)+O(|g_{\mathrm Q}|^4),\quad
 P_0=1-|g_{\mathrm Q}|^2(2\bar n+1)+O(|g_{\mathrm Q}|^4).
 \label{eq:feqo-weak-quantum-populations}
\end{equation}
此外
\begin{equation}
 (\rho_e)_{1,0}=g_{\mathrm Q}\langle a\rangle+O(|g_{\mathrm Q}|^2),\qquad
 (\rho_e)_{-1,0}=-g_{\mathrm Q}^*\langle a^\dagger\rangle+O(|g_{\mathrm Q}|^2).
 \label{eq:feqo-quantum-coherence-first-order}
\end{equation}
\end{proposition}
\begin{proof}
吸收末态光向量为 $g_{\mathrm Q}a|\gamma\rangle$，取范数平方给
$|g_{\mathrm Q}|^2\langle a^\dagger a\rangle$；发射对应
$|g_{\mathrm Q}|^2\langle aa^\dagger\rangle$。二阶返回项保证总概率为一。
吸收向量与无交换光向量的重叠给 $g_{\mathrm Q}\langle a\rangle$，发射同理。
混合输入由线性性成立。有限阶余项要求相应光子数矩有限，且显著占据扇区满足
$|g_{\mathrm Q}|\sqrt{n+1}\ll1$；有限失谐则须使用逐边的时间积分。
\end{proof}
平均光子数相同的 Fock、热态与相干态在最低阶人口上可以相同，
但前两者 $\langle a\rangle=0$，相干态保留一阶电子相干。
高阶增益的领先概率为
$P_{+k}=|g_{\mathrm Q}|^{2k}\langle(N)_k\rangle/(k!)^2$。
固定 $\bar n$ 后，热态 $\langle(N)_k\rangle=k!\bar n^k$，相干态为 $\bar n^k$，
因此两种增益概率的比值为 $k!$；同时绝对概率仍受 $(k!)^{-2}$ 抑制。
\begin{figure}[htbp]
 \centering
 \includegraphics[width=0.76\linewidth]{photon_statistics_gain.pdf}
 \caption{相同 $\bar n=3$ 下的高阶增益系数 $P_{+k}/|g_{\mathrm Q}|^{2k}$。
 相干态为 $\bar n^k/(k!)^2$，单模热态为 $\bar n^k/k!$；两者比值为 $k!$。
 图示领先微扰系数，不能把任意大的 $k$ 系数直接当成有限耦合概率。}
 \label{fig:photon-statistics-gain}
\end{figure}
